We ask whether lightweight vision--language models (VLMs) can be adapted to
medical visual question answering (VQA) on consumer-grade GPUs, and answer it
with a three-stage study under a single evaluation protocol. Four open 2--3B
models (Qwen3-VL-2B, Qwen2.5-VL-3B, SmolVLM2-2.2B, Gemma4-E2B) are evaluated on
PathVQA, SLAKE and VQA-RAD across zero-shot baselines (12 conditions), QLoRA
fine-tuning (75 conditions), and a comparison of hyperparameter search
strategies (810 trials). Zero-shot accuracy differs significantly across models
(Cochran's $Q = 1904.28$, $\mathit{df} = 3$, $p < .001$), and the ranking
survives removal of samples flagged as likely pretraining exposure by
Min-$K$\% Probability. The effect of QLoRA is heterogeneous \emph{in direction}:
it significantly improves Qwen2.5-VL-3B ($d = +2.646$) and Qwen3-VL-2B
($d = +1.620$) yet significantly degrades SmolVLM2-2.2B ($d = -2.284$), and a
mixed-effects model pooling all four reports no effect ($p = .3629$) because the
opposing effects cancel. In autonomous hyperparameter search, an LLM agent
($0.4184$) performs significantly worse than Bayesian optimization ($0.4490$;
Mann--Whitney $U = 16.00$, $p = .0112$, $r = -0.68$) and is indistinguishable
from random search ($0.4186$), though all three automated strategies exceed the
manual configuration ($0.3776$). The agent's disadvantage traces not to proposal
quality but to a reduced effective budget caused by duplicate proposals ($12.5$
of $20$ unique configurations per repeat); a 200-trial re-experiment removing
the prompt--implementation inconsistency recovered uniqueness to $18.90$ of $20$
and dissolved the significant gap ($p = .1726$), yet still did not demonstrate
improvement ($p = .2387$). Notably, the variation observed when repeating an
identical configuration ($0.056$) exceeds the mean difference between strategies
($0.031$), a quantitative counterexample to comparing search methods from single
runs.
